% 2007 年湖北卷文科第 6 题：100 名男生体重的频率分布直方图。
%   组距 2 kg，横轴 52.5—76.5 kg，纵轴「频率/组距」0.01—0.08。
% 来源 PDF：gaokaomath-src/普通高考/2007/2007湖北文.pdf
% SHA-256：290a153da24ec0f3a3abb8bb916d9b7e06d203d3ef5bd7bb221c93f7f1b1cacb
% 原页：2007湖北文.pdf 第 1 页（A3 三栏）第 1 栏，书页码 377。
% 重绘依据：按原页 12 个等宽柱与纵轴刻度重绘，柱高按原图读取。
\begin{tikzpicture}[x=1.0cm, y=70cm, semithick,
  every node/.style={inner sep=1pt, font=\scriptsize}]
  \draw[->,>=stealth] (0,0) -- (12.8,0) node[right] {体重 (kg)};
  \draw[->,>=stealth] (0,0) -- (0,0.09) node[above] {频率/组距};

  \foreach \i/\h in {0/0.01,1/0.01,2/0.02,3/0.05,4/0.06,5/0.07,
                     6/0.08,7/0.065,8/0.05,9/0.04,10/0.03,11/0.02} {
    \draw (\i,0) rectangle (\i+1,\h);
  }

  \foreach \y in {0.01,0.02,0.03,0.04,0.05,0.06,0.07,0.08} {
    \draw (0,\y) -- (-0.12,\y);
    \node[left] at (0,\y) {\y};
  }
  \node[below left] at (0,0) {$O$};
  \foreach \i in {1,...,12} {
    \draw (\i,0) -- (\i,-0.0015);
    \node[below] at (\i,0) {\the\numexpr 52+2*\i\relax.5};
  }
\end{tikzpicture}
